\documentclass[12pt,a4paper]{article}
\usepackage[utf8]{inputenc}
\usepackage{amsmath}
\usepackage{amsfonts}
\usepackage{amssymb}
\usepackage{graphicx}
\usepackage{geometry}
\usepackage{fancyhdr}
\usepackage{hyperref}

\geometry{top=1in,bottom=1in,left=1in,right=1in}

\pagestyle{fancy}
\fancyhf{}
\rhead{FORM 2 -- COMPLETE SPECIFICATION}
\lfoot{APPLICANT: JASON PANDIAN}
\rfoot{Page \thepage}
\renewcommand{\headrulewidth}{0.4pt}
\renewcommand{\footrulewidth}{0.4pt}

\begin{document}

\begin{center}
    \textbf{\large FORM 2} \\
    \vspace{0.3cm}
    \textbf{THE PATENTS ACT, 1970} \\
    (39 of 1970) \\
    \& \\
    \textbf{THE PATENT RULES, 2003} \\
    \vspace{0.5cm}
    \textbf{\large COMPLETE SPECIFICATION} \\
    (See Section 10 and Rule 13) \\
    \vspace{1cm}
    
    \textbf{\Large 1. TITLE OF THE INVENTION} \\
    \vspace{0.5cm}
    \textbf{\large A PREDICTIVE TEMPORAL BRIDGE ARCHITECTURE WITH DEDICATED ARTIFICIAL INTELLIGENCE HARDWARE FOR VIRTUAL REAL-TIME DEEP SPACE COMMUNICATION} \\
    \vspace{1.5cm}
    
    \textbf{\Large 2. APPLICANT(S)} \\
    \vspace{0.5cm}
    \begin{tabular}{lll}
        \textbf{Name} & : & Jason Pandian \\
        \textbf{Nationality} & : & Indian National \\
        \textbf{Address} & : & Nehru Institute of Technology, \\
                         &   & Coimbatore, Tamilnadu, India. \\
    \end{tabular}
    \vspace{1.5cm}
    
    \textbf{\Large 3. PREAMBLE TO THE DESCRIPTION} \\
    \vspace{0.5cm}
    \textit{The following specification particularly describes the invention and the manner in which it is to be performed:}
\end{center}

\newpage

\section*{FIELD OF THE INVENTION}
The present disclosure generally relates to the field of deep-space communications and artificial intelligence. More particularly, the present invention relates to a Predictive Temporal Bridge (PTB) architecture and a dedicated, low-power, radiation-hardened AI co-processor chip designed to fundamentally mask extreme communication light-time delays, enabling virtual real-time interaction between ground control stations and interplanetary spacecraft while maintaining human autonomy and system safety.

\section*{BACKGROUND OF THE INVENTION AND PRIOR ART}
Interplanetary spaceflight, particularly crewed missions to Mars, faces an insurmountable physical constraint: the speed of light. At median Mars-Earth opposition, radio and optical signals take approximately 4 to 22 minutes to travel one-way. This communication latency creates a severe temporal disconnection between Ground Control and the spacecraft crew.

Traditionally, deep-space operations operate on asynchronous telemetry. Ground controllers view stale data (from the past), formulate commands, and transmit them back, resulting in a round-trip delay that shatters real-time collaboration. This induces a high Cognitive Load Index (CLI), lowers the Social Presence Index (SPI), and destroys Shared Situational Awareness (SA). During mission-critical anomalies, this latency can be fatal, as ground support cannot provide immediate intervention.

Prior art attempts to solve this include:
\begin{enumerate}
    \item \textbf{High-Degree Spacecraft Autonomy:} Pushing decision-making entirely to the spacecraft. While effective for safety, it isolates the crew psychologically and limits their access to Earth's massive computational resources.
    \item \textbf{Predictive Mechanical Teleoperation (e.g., Smith Predictors):} Used for robotic arms by visually predicting the state of the arm. However, this is limited to single-operator mechanical tasks and does not synchronize cognitive and social interaction frames between two separated human teams.
\end{enumerate}

To process physics-based predictive models locally on a spacecraft, complex Artificial Intelligence (AI) software is typically required. Running these models on standard spacecraft central processing units (CPUs) or field-programmable gate arrays (FPGAs) draws significant power, lacks deterministic execution speeds required for real-time reconciliation, and competes with critical flight software for CPU cycles. 

Therefore, there is an unaddressed need for a comprehensive system that can bridge the temporal gap using synchronized AI models, coupled with a dedicated, special-purpose AI hardware chip that requires minimal power, offloads processing from the main flight computer, and interfaces seamlessly with existing spacecraft avionics and ground station architectures.

\section*{OBJECT OF THE INVENTION}
The primary object of the present invention is to provide a Predictive Temporal Bridge (PTB) architecture that masks deep-space communication latency by mathematically projecting ground intent forward in time, establishing a virtual synchronous environment.

Another object is to provide a dedicated, specialized, low-power AI co-processor chip (the Predictive Temporal Processing Unit, or PTPU) in which predictive physics and reconciliation logic are preprogrammed as unalterable, deterministic firmware.

Yet another object of the invention is to design the PTPU chip to be universally adaptable, featuring standard aerospace bus interfaces (such as SpaceWire or MIL-STD-1553) allowing it to be easily interfaced with any existing ground station terminal and spacecraft flight computer.

A further object is to provide a dual-gated reality reconciliation method executed in physical silicon to ensure that AI-predicted ground commands are validated against local sensor truth prior to human authorization, preventing automation bias.

\section*{SUMMARY OF THE INVENTION}
The present invention achieves the aforementioned objectives by providing a Predictive Temporal Bridge (PTB) comprising a Ground Prediction Module and a Spacecraft Reconciliation Module, anchored by specialized hardware. 

The core of the invention lies in the dedicated hardware: the Predictive Temporal Processing Unit (PTPU). The PTPU is a low-power, radiation-hardened Application-Specific Integrated Circuit (ASIC) designed as a "plug-and-play" co-processor for existing space and ground systems. The PTPU encapsulates a minimal, deterministic CPU architecture (such as a dual-core lockstep RISC-V or ARM Cortex-R) specifically tuned to execute physics-based machine learning inference and matrix calculations without interrupting the main flight computer. The PTB algorithms are burnt into the PTPU's read-only non-volatile memory (firmware), ensuring absolute execution reliability.

In operation, a Ground Station equipped with a PTPU receives asynchronous telemetry from the spacecraft (delayed by time $m$). The Ground PTPU rapidly computes a forward-projection of the spacecraft's state to time $t+2m$ and overlays proposed correction commands. This preemptive command bundle is transmitted to the spacecraft. 

When the command arrives at the spacecraft, the Spacecraft PTPU intercepts the packet via its aerospace bus interface. It executes a low-latency "reality overlay" algorithm, directly comparing the Earth AI's predicted state against the spacecraft's real-time local sensors. If the divergence between prediction and reality is within safety margins, the PTPU flags the command as validated for the human commander to execute, reducing perceived latency to near-zero. If unmodeled physical anomalies exist, the PTPU instantly blocks execution, preserving crew autonomy.

\section*{BRIEF DESCRIPTION OF THE ACCOMPANYING DRAWINGS}
\begin{description}
    \item[FIG. 1] is a high-level system architecture diagram of the Predictive Temporal Bridge (PTB) linking a Ground Station and a Deep-Space Spacecraft, illustrating the temporal projection logic.
    \item[FIG. 2] is a detailed hardware block diagram of the Predictive Temporal Processing Unit (PTPU) AI chip, showing its internal minimal CPU, matrix accelerators, firmware logic, and aerospace bus interfaces.
    \item[FIG. 3] is a state transition diagram of the Reality Reconciliation Engine executing within the Spacecraft PTPU.
    \item[FIG. 4] illustrates the integration of the PTPU chip into an existing spacecraft avionics network.
\end{description}

\section*{DETAILED DESCRIPTION OF THE INVENTION}

\subsection*{1. System Architecture of the Predictive Temporal Bridge (FIG. 1)}
The overall PTB architecture establishes a virtual real-time communication channel. A Deep-Space Spacecraft operating at a physical time $t$ generates a telemetry broadcast packet comprising state vectors, sensor noise envelopes, and human intent vectors. This packet travels across space, arriving at a Ground Station at time $t+m$, where $m$ is the one-way light-time delay (e.g., 20 minutes).

At the Ground Station, the telemetry is fed into a Ground PTPU chip. Instead of displaying only the delayed state, the Ground PTPU utilizes physics-based AI to project the state forward to $t+2m$, estimating the spacecraft's condition at the exact moment a response will arrive. Ground Controllers formulate a correction command based on this projected future state, and transmit it back to the spacecraft.

Upon arrival at the spacecraft at time $t+2m$, the preemptive command is captured and processed by a Spacecraft PTPU chip. The Spacecraft PTPU acts as a Reality Arbiter, instantaneously comparing the Earth's predictive model with the local, present physical reality before presenting the command to the spacecraft commander on an interactive dashboard. 

\subsection*{2. Hardware Design of the PTPU AI Chip (FIG. 2 \& FIG. 4)}
To guarantee deterministic execution and minimal power overhead, the PTB logic is not run as standard software on the main flight computer. Instead, it is offloaded to the Predictive Temporal Processing Unit (PTPU).

The PTPU is a specialized, low-power (e.g., < 2 Watts) AI co-processor chip. The internal hardware specifications comprise:
\begin{itemize}
    \item \textbf{Minimal Deterministic CPU Core:} A highly reliable, dual-core lockstep processor (such as a rad-hard RISC-V or ARM Cortex-R5) functioning as the orchestrator. It is devoid of heavy operating systems, running bare-metal logic to ensure microsecond-level determinism.
    \item \textbf{Neural Matrix Accelerator (NMA):} A dedicated hardware block of multiply-accumulate (MAC) units specifically optimized for evaluating pre-trained kinematic and thermodynamic neural network models.
    \item \textbf{Non-Volatile Firmware ROM:} The entire PTB prediction, interpolation, and reconciliation algorithms are preprogrammed into physical silicon ROM. This eliminates the risk of software bit-flips in deep space and prevents unauthorized code modification.
    \item \textbf{Universal Aerospace Bus Interface Controller:} A built-in hardware controller featuring native support for SpaceWire, MIL-STD-1553, and CAN bus protocols. This ensures the chip can be seamlessly integrated into legacy ground station racks or connected directly to any standard spacecraft avionics bus.
    \item \textbf{Local Fast SRAM:} Error-correcting code (ECC) static RAM for temporarily storing incoming telemetry packets and intermediate matrix computations.
\end{itemize}

As shown in FIG. 4, the PTPU chip acts as a peripheral node on the spacecraft's primary SpaceWire bus. It receives telemetry from sensors in parallel with the main Flight Computer, performs its AI reconciliation tasks autonomously, and outputs validated "safe-to-execute" command flags to the Human-Machine Interface (HMI) displays, offloading entirely the heavy predictive processing from the main flight computer.

\subsection*{3. Reality Reconciliation Engine (FIG. 3)}
The PTPU executes a hardware-level Finite State Machine (FSM) for reality reconciliation. When a predicted command packet arrives from Earth:
\begin{enumerate}
    \item \textbf{State 1 (Data Ingress):} The PTPU extracts the Earth AI's predicted state vector $\hat{T}(t_{current})$.
    \item \textbf{State 2 (Sensor Sampling):} The PTPU reads the immediate physical sensor telemetry $T(t_{current})$.
    \item \textbf{State 3 (Deviation Calculation):} The minimal CPU calculates the variance matrix $\Delta = |T(t_{current}) - \hat{T}(t_{current})|$.
    \item \textbf{State 4 (Threshold Gating):} The variance $\Delta$ is compared against a hardcoded safety tolerance $\delta_{safe}$.
    \item \textbf{State 5 (Reconciliation/Flagging):} If $\Delta < \delta_{safe}$, the FSM transitions to the "Reconciled" state, sending a green authorization token to the commander's dashboard. If $\Delta > \delta_{safe}$, it transitions to the "Desynced" state, blocking auto-execution and demanding a manual override due to an unmodeled local anomaly.
\end{enumerate}

By executing this logic entirely within the dedicated hardware, the system ensures that the cognitive load on the human operator is minimized, while preventing automation bias by strictly enforcing mathematical verification of AI predictions.

\newpage

\section*{WE CLAIM:}

\begin{enumerate}
    \item A Predictive Temporal Bridge (PTB) system for mitigating communication latency in deep-space networks, comprising:
    \begin{itemize}
        \item a Ground Predictive Temporal Processing Unit (PTPU) deployed at an Earth ground station, configured to receive asynchronous spacecraft telemetry and mathematically project said telemetry forward in time to generate a predicted future state vector;
        \item a Spacecraft Predictive Temporal Processing Unit (PTPU) deployed on an interplanetary spacecraft, configured to receive preemptive commands based on the predicted future state vector; and
        \item a hardware-based Reality Reconciliation Engine executed within the Spacecraft PTPU, configured to instantly compare the received predicted future state vector against real-time local sensor telemetry to generate an execution authorization flag.
    \end{itemize}

    \item The system of claim 1, wherein the Ground PTPU and the Spacecraft PTPU are identical dedicated AI co-processor chips, operating as low-power, radiation-hardened Application-Specific Integrated Circuits (ASICs).

    \item The system of claim 2, wherein the PTPU chip comprises:
    \begin{itemize}
        \item a minimal, deterministic dual-core lockstep central processing unit (CPU);
        \item a dedicated Neural Matrix Accelerator (NMA) hardware block optimized for continuous kinematic and thermodynamic neural network inference;
        \item a Non-Volatile Firmware memory wherein predictive modeling and reality reconciliation algorithms are preprogrammed as unalterable logic; and
        \item a Universal Aerospace Bus Interface Controller configured for direct physical connection to standard spacecraft data buses including SpaceWire, MIL-STD-1553, and CAN bus.
    \end{itemize}

    \item The system of claim 3, wherein the PTPU chip consumes less than 2 Watts of continuous power and operates independently of the spacecraft's main flight computer, functioning as a peripheral AI co-processor.

    \item A method for achieving virtual real-time interaction in extreme latency environments using a dedicated AI hardware co-processor, the method comprising:
    \begin{itemize}
        \item capturing, by a Ground PTPU, a telemetry packet generated by a remote spacecraft at time $t$ and received at time $t+m$, where $m$ is the signal propagation delay;
        \item executing, within the Ground PTPU's neural matrix accelerator, a forward-projection algorithm to output a predicted state vector at time $t+2m$;
        \item formulating a ground command packet encapsulated with the predicted state vector, and transmitting the packet to the remote spacecraft;
        \item intercepting the packet, via an aerospace bus interface, by a Spacecraft PTPU upon arrival at time $t+2m$;
        \item calculating, using a minimal deterministic CPU within the Spacecraft PTPU, a mathematical variance between the predicted state vector and real-time physical sensor data;
        \item validating the ground command if the variance is below a predefined hardware tolerance, and flagging it for human execution; and
        \item blocking auto-execution of the ground command if the variance exceeds the predefined hardware tolerance, thereby enforcing human autonomy during unmodeled physical anomalies.
    \end{itemize}

    \item The method of claim 5, wherein the validation process executed by the Spacecraft PTPU completes in deterministic microsecond timing without utilizing processing cycles from the spacecraft's primary flight software.
\end{enumerate}

\newpage

\section*{ABSTRACT OF THE DISCLOSURE}
\textbf{A PREDICTIVE TEMPORAL BRIDGE ARCHITECTURE WITH DEDICATED ARTIFICIAL INTELLIGENCE HARDWARE FOR VIRTUAL REAL-TIME DEEP SPACE COMMUNICATION} \\
\vspace{0.5cm}

A Predictive Temporal Bridge (PTB) architecture and a dedicated artificial intelligence hardware co-processor are disclosed for eliminating the operational and psychological impacts of severe light-time communication delays in deep space. The system comprises a pair of identical, low-power Predictive Temporal Processing Unit (PTPU) chips installed at both Ground Control and the Spacecraft. The PTPU is a rad-hardened ASIC containing a deterministic minimal CPU, a neural matrix accelerator, and universal aerospace bus interfaces (e.g., SpaceWire). The Ground PTPU mathematically projects incoming delayed telemetry forward by the round-trip latency time, allowing controllers to formulate preemptive commands. Upon receipt, the Spacecraft PTPU executes a hardware-level reality reconciliation engine, immediately validating the AI-predicted ground command against real-time local sensor truth before human authorization. This creates a virtual real-time interactive environment while preserving system safety and human autonomy.

\end{document}
